\chapter{Thời gian và không gian}
\markboth{Thời gian và không gian}{Time and Space}


\section{Không ai có thể nghĩ hộ em — im lặng là lúc em nghĩ cho mình.}
\begin{truyen}{Không ai có thể nghĩ hộ em — im lặng là lúc em nghĩ cho mình...}{Không ai có thể nghĩ hộ em — im lặng là lúc em ngh}
\chuhoa{C}{âu thầy nói câu ấy. Cả lớp im lặng — không ai nói, mỗi người đang nghĩ về điều vừa nghe.}
\textit{(The teacher said that. The class fell silent — no one spoke; everyone was thinking about what they had just heard.)}
\end{truyen}

\begin{baihoc}
Im lặng sau câu nói đúng lúc là dấu hiệu của tư duy.
\textit{(Silence after the right sentence is a sign of thinking.)}
\end{baihoc}

\begin{ghinhoanh}
\textbf{Short English to remember:}

Some sentences make the whole class think in silence.

\end{ghinhoanh}

\begin{tuhocnhanh}
1. Câu này khiến em nghĩ đến điều gì? \textit{What does this make you think of?}

2. Em có thể dùng câu này trong tình huống nào? \textit{When could you use this sentence?}
\end{tuhocnhanh}
\ngancach


\section{Em có quyền im lặng — và quyền đó rất quý.}
\begin{truyen}{Em có quyền im lặng — và quyền đó rất quý.}{Em có quyền im lặng — và quyền đó rất quý.}
\chuhoa{C}{âu thầy nói câu ấy. Cả lớp im lặng — không ai nói, mỗi người đang nghĩ về điều vừa nghe.}
\textit{(The teacher said that. The class fell silent — no one spoke; everyone was thinking about what they had just heard.)}
\end{truyen}

\begin{baihoc}
Im lặng sau câu nói đúng lúc là dấu hiệu của tư duy.
\textit{(Silence after the right sentence is a sign of thinking.)}
\end{baihoc}

\begin{ghinhoanh}
\textbf{Short English to remember:}

Some sentences make the whole class think in silence.

\end{ghinhoanh}

\begin{tuhocnhanh}
1. Câu này khiến em nghĩ đến điều gì? \textit{What does this make you think of?}

2. Em có thể dùng câu này trong tình huống nào? \textit{When could you use this sentence?}
\end{tuhocnhanh}
\ngancach


\section{Nói ít, nghĩ nhiều — đó là cách người lớn dùng lời.}
\begin{truyen}{Nói ít, nghĩ nhiều — đó là cách người lớn dùng lời.}{Nói ít, nghĩ nhiều — đó là cách người lớn dùng lời}
\chuhoa{C}{âu thầy nói câu ấy. Cả lớp im lặng — không ai nói, mỗi người đang nghĩ về điều vừa nghe.}
\textit{(The teacher said that. The class fell silent — no one spoke; everyone was thinking about what they had just heard.)}
\end{truyen}

\begin{baihoc}
Im lặng sau câu nói đúng lúc là dấu hiệu của tư duy.
\textit{(Silence after the right sentence is a sign of thinking.)}
\end{baihoc}

\begin{ghinhoanh}
\textbf{Short English to remember:}

Some sentences make the whole class think in silence.

\end{ghinhoanh}

\begin{tuhocnhanh}
1. Câu này khiến em nghĩ đến điều gì? \textit{What does this make you think of?}

2. Em có thể dùng câu này trong tình huống nào? \textit{When could you use this sentence?}
\end{tuhocnhanh}
\ngancach


\section{Khi cả lớp im, có thể đang có điều quan trọng được suy nghĩ.}
\begin{truyen}{Khi cả lớp im, có thể đang có điều quan trọng được suy nghĩ.}{Khi cả lớp im, có thể đang có điều quan trọng được}
\chuhoa{C}{âu thầy nói câu ấy. Cả lớp im lặng — không ai nói, mỗi người đang nghĩ về điều vừa nghe.}
\textit{(The teacher said that. The class fell silent — no one spoke; everyone was thinking about what they had just heard.)}
\end{truyen}

\begin{baihoc}
Im lặng sau câu nói đúng lúc là dấu hiệu của tư duy.
\textit{(Silence after the right sentence is a sign of thinking.)}
\end{baihoc}

\begin{ghinhoanh}
\textbf{Short English to remember:}

Some sentences make the whole class think in silence.

\end{ghinhoanh}

\begin{tuhocnhanh}
1. Câu này khiến em nghĩ đến điều gì? \textit{What does this make you think of?}

2. Em có thể dùng câu này trong tình huống nào? \textit{When could you use this sentence?}
\end{tuhocnhanh}
\ngancach


\section{Điều em giữ trong im lặng đôi khi mạnh hơn điều em nói ra.}
\begin{truyen}{Điều em giữ trong im lặng đôi khi mạnh hơn điều em nói ra.}{Điều em giữ trong im lặng đôi khi mạnh hơn điều em}
\chuhoa{C}{âu thầy nói câu ấy. Cả lớp im lặng — không ai nói, mỗi người đang nghĩ về điều vừa nghe.}
\textit{(The teacher said that. The class fell silent — no one spoke; everyone was thinking about what they had just heard.)}
\end{truyen}

\begin{baihoc}
Im lặng sau câu nói đúng lúc là dấu hiệu của tư duy.
\textit{(Silence after the right sentence is a sign of thinking.)}
\end{baihoc}

\begin{ghinhoanh}
\textbf{Short English to remember:}

Some sentences make the whole class think in silence.

\end{ghinhoanh}

\begin{tuhocnhanh}
1. Câu này khiến em nghĩ đến điều gì? \textit{What does this make you think of?}

2. Em có thể dùng câu này trong tình huống nào? \textit{When could you use this sentence?}
\end{tuhocnhanh}
\ngancach


\section{Học cách im lặng cũng là học.}
\begin{truyen}{Học cách im lặng cũng là học.}{Học cách im lặng cũng là học.}
\chuhoa{C}{âu thầy nói câu ấy. Cả lớp im lặng — không ai nói, mỗi người đang nghĩ về điều vừa nghe.}
\textit{(The teacher said that. The class fell silent — no one spoke; everyone was thinking about what they had just heard.)}
\end{truyen}

\begin{baihoc}
Im lặng sau câu nói đúng lúc là dấu hiệu của tư duy.
\textit{(Silence after the right sentence is a sign of thinking.)}
\end{baihoc}

\begin{ghinhoanh}
\textbf{Short English to remember:}

Some sentences make the whole class think in silence.

\end{ghinhoanh}

\begin{tuhocnhanh}
1. Câu này khiến em nghĩ đến điều gì? \textit{What does this make you think of?}

2. Em có thể dùng câu này trong tình huống nào? \textit{When could you use this sentence?}
\end{tuhocnhanh}
\ngancach


\section{Một khoảnh khắc im lặng có thể đáng giá cả trang sách.}
\begin{truyen}{Một khoảnh khắc im lặng có thể đáng giá cả trang sách.}{Một khoảnh khắc im lặng có thể đáng giá cả trang s}
\chuhoa{C}{âu thầy nói câu ấy. Cả lớp im lặng — không ai nói, mỗi người đang nghĩ về điều vừa nghe.}
\textit{(The teacher said that. The class fell silent — no one spoke; everyone was thinking about what they had just heard.)}
\end{truyen}

\begin{baihoc}
Im lặng sau câu nói đúng lúc là dấu hiệu của tư duy.
\textit{(Silence after the right sentence is a sign of thinking.)}
\end{baihoc}

\begin{ghinhoanh}
\textbf{Short English to remember:}

Some sentences make the whole class think in silence.

\end{ghinhoanh}

\begin{tuhocnhanh}
1. Câu này khiến em nghĩ đến điều gì? \textit{What does this make you think of?}

2. Em có thể dùng câu này trong tình huống nào? \textit{When could you use this sentence?}
\end{tuhocnhanh}
\ngancach


\section{Thầy chờ không phải vì thầy không biết — thầy chờ để em có thời gian nghĩ.}
\begin{truyen}{Thầy chờ không phải vì thầy không biết — thầy chờ để em có t...}{Thầy chờ không phải vì thầy không biết — thầy chờ }
\chuhoa{C}{âu thầy nói câu ấy. Cả lớp im lặng — không ai nói, mỗi người đang nghĩ về điều vừa nghe.}
\textit{(The teacher said that. The class fell silent — no one spoke; everyone was thinking about what they had just heard.)}
\end{truyen}

\begin{baihoc}
Im lặng sau câu nói đúng lúc là dấu hiệu của tư duy.
\textit{(Silence after the right sentence is a sign of thinking.)}
\end{baihoc}

\begin{ghinhoanh}
\textbf{Short English to remember:}

Some sentences make the whole class think in silence.

\end{ghinhoanh}

\begin{tuhocnhanh}
1. Câu này khiến em nghĩ đến điều gì? \textit{What does this make you think of?}

2. Em có thể dùng câu này trong tình huống nào? \textit{When could you use this sentence?}
\end{tuhocnhanh}
\ngancach


\section{Câu nói khiến em im lặng là câu nói đáng nhớ.}
\begin{truyen}{Câu nói khiến em im lặng là câu nói đáng nhớ.}{Câu nói khiến em im lặng là câu nói đáng nhớ.}
\chuhoa{C}{âu thầy nói câu ấy. Cả lớp im lặng — không ai nói, mỗi người đang nghĩ về điều vừa nghe.}
\textit{(The teacher said that. The class fell silent — no one spoke; everyone was thinking about what they had just heard.)}
\end{truyen}

\begin{baihoc}
Im lặng sau câu nói đúng lúc là dấu hiệu của tư duy.
\textit{(Silence after the right sentence is a sign of thinking.)}
\end{baihoc}

\begin{ghinhoanh}
\textbf{Short English to remember:}

Some sentences make the whole class think in silence.

\end{ghinhoanh}

\begin{tuhocnhanh}
1. Câu này khiến em nghĩ đến điều gì? \textit{What does this make you think of?}

2. Em có thể dùng câu này trong tình huống nào? \textit{When could you use this sentence?}
\end{tuhocnhanh}
\ngancach


\section{Im lặng sau câu hỏi là dấu hiệu lớp đang làm việc.}
\begin{truyen}{Im lặng sau câu hỏi là dấu hiệu lớp đang làm việc.}{Im lặng sau câu hỏi là dấu hiệu lớp đang làm việc.}
\chuhoa{C}{âu thầy nói câu ấy. Cả lớp im lặng — không ai nói, mỗi người đang nghĩ về điều vừa nghe.}
\textit{(The teacher said that. The class fell silent — no one spoke; everyone was thinking about what they had just heard.)}
\end{truyen}

\begin{baihoc}
Im lặng sau câu nói đúng lúc là dấu hiệu của tư duy.
\textit{(Silence after the right sentence is a sign of thinking.)}
\end{baihoc}

\begin{ghinhoanh}
\textbf{Short English to remember:}

Some sentences make the whole class think in silence.

\end{ghinhoanh}

\begin{tuhocnhanh}
1. Câu này khiến em nghĩ đến điều gì? \textit{What does this make you think of?}

2. Em có thể dùng câu này trong tình huống nào? \textit{When could you use this sentence?}
\end{tuhocnhanh}
\ngancach
